\section{Introduction}
\label{sec:introduction}

We study how many gates are needed to compute iterated matrix multiplication
with five homogeneous arithmetic layers.  For integers $w,d\geq 1$, let
$X^{(1)},\ldots,X^{(d)}$ be $w\times w$ matrices whose entries are distinct
commuting variables.  Define
\begin{equation}
\label{eq:imm-definition}
 \IMM_{w,d}
 =\bigl(X^{(1)}\cdots X^{(d)}\bigr)_{1,1}
 =\sum_{i_1,\ldots,i_{d-1}\in[w]}
   x^{(1)}_{1,i_1}x^{(2)}_{i_1,i_2}\cdots x^{(d)}_{i_{d-1},1}.
\end{equation}
For $d=1$, the expression means $x^{(1)}_{1,1}$.  The first index denotes
matrix width and the second denotes degree.  We retain all $dw^2$ matrix-entry
variables as ambient variables, including entries absent from the polynomial.
Our target is $\IMM_{n,n}$, of degree $n$ in $n^3$ ambient variables.

\subsection{Circuit model and main result}

An arithmetic circuit over a field $K$ is a finite directed acyclic graph.
Its leaves are variables or elements of $K$.  A sum gate computes a
$K$-linear combination of its inputs, and a product gate computes their
product.  Fan-in and fan-out are arbitrary and finite.  Gates may be shared,
and a gate may occur repeatedly among another gate's inputs; repeated
occurrences at a product gate are counted with their multiplicities.
Scalar coefficients on sum inputs are part of the circuit and do not require
additional gates.

A $\Sigma\Pi\Sigma\Pi\Sigma$ circuit has five computational layers,
listed from output to leaves, with respective gate types
$+,\times,+,\times,+$ and a single output gate.  Edges between computational
gates join consecutive layers, and the bottom sums read leaves.  Unary gates
are allowed.  In particular, no bound is placed on the number of variables
in a bottom linear form.

The circuit is \emph{syntactically homogeneous} if it has the following
formal-degree assignment: variables have degree $1$, constants have degree
$0$, product gates add their input degrees with multiplicity, and all inputs
of a sum gate have the same degree, which is assigned to that gate.  Empty
sums and products, if present, compute $0$ and $1$ and have formal degree $0$.
Thus a zero polynomial at a gate still carries its assigned formal degree;
cancellation does not reset that degree.  All computations and equalities
below concern formal polynomials.

We define \emph{size} to be the number of vertices, including leaves.  The
number of computational gates excludes leaves, whereas wire size counts
input occurrences.  These quantities are not identified: our lower bound is
on gates, and the circuit analysis also bounds the number of computational
gates alone.

\begin{theorem}
\label{thm:main}
There exist absolute constants $c>0$ and $n_0$ such that, for every integer
$n\geq n_0$, every syntactically homogeneous
$\Sigma\Pi\Sigma\Pi\Sigma$ circuit over $\C$ computing $\IMM_{n,n}$
has size at least
\[
 n^{c\sqrt n}.
\]
One may take $c=1/400$.
\end{theorem}

The model permits unrestricted bottom fan-in and support, constants, and
shared subcircuits.  Corollary~\ref{cor:characteristic-zero} extends the
$n^{\sqrt n/400}$ lower bound to every field of characteristic zero, with
the same absolute threshold.  Proposition~\ref{prop:upper} gives an
elementary block construction of size at most $n^{\sqrt n+4}$ over every
field for $n\geq2$.  It is homogeneous of depth four; inserting unary bottom
sums makes it depth five, so the lower bound has the matching scale.
Iterated matrix multiplication nevertheless has polynomial-size circuits
when depth is unrestricted.

\subsection{Historical context and previous bounds}

Iterated matrix multiplication has long served as a test of the cost of
restricting arithmetic depth.  Its path expansion connects matrix products
with reachability, while its polynomial-size unrestricted circuits make it
a natural target for lower bounds on weaker circuit models.
Nisan and Wigderson used dimensions of partial-derivative spaces to prove
lower bounds for homogeneous depth-three circuits and for constant-depth
set-multilinear circuits.  They asked for an explicit homogeneous polynomial
requiring superpolynomial size at constant depth, even at depth five
\cite[Section~2.2 of the linked author journal manuscript]{NisanWigderson1995}.

Depth reduction gave this program a further motivation.
Agrawal and Vinay reduced general arithmetic computation to depth four,
and Koiran and Tavenas sharpened the size bounds
\cite{AgrawalVinay2008,Koiran2012,Tavenas2013}.
For $d=N^{O(1)}$, Tavenas's theorem converts a circuit of size
$N^{O(1)}$ computing a homogeneous degree-$d$ polynomial in $N$ variables
to a homogeneous depth-four circuit of size $N^{O(\sqrt d)}$
\cite[Theorem~1 of the linked preprint]{Tavenas2013}.
This explains the square-root scale in depth-four lower bounds and in the
block construction for IMM used here.

The progress toward this scale involved several distinct restrictions.
Fournier, Limaye, Malod, and Srinivasan proved
IMM lower bounds with bounds on both product-layer fan-ins and, separately,
for homogeneous multilinear formulas
\cite[Theorems~16 and~29 of the cited author version]{FournierLimayeMalodSrinivasan2015}.
Kayal, Limaye, Saha, and Srinivasan removed the bottom-fan-in restriction
for homogeneous depth-four formulas, obtaining an exponential lower bound
over characteristic zero for a Nisan--Wigderson design family
\cite[Theorem~1]{KayalLimayeSahaSrinivasan2014}.
Kumar and Saraf proved a $d^{\Omega(\sqrt d)}$ lower bound for homogeneous
depth-four circuits computing $\IMM_{d^5,d}$, over every field
\cite[Theorem~1.2 and Section~8.2 of the cited preprint]{KumarSaraf2014}.
In another direction, Kayal, Saha, and Tavenas obtained
$w^{\Omega(\sqrt d)}$ for syntactically multilinear depth-four circuits
computing $\IMM_{w,d}$, without homogeneity, in the range
$d\geq\log^2 w$, over every field
\cite[Theorem~1.4 and the following comparison]{KayalSahaTavenas2018}.

At depth five, Bera and Chakrabarti proved a lower bound with a restriction
on bottom support.
Their Theorem~1.1 states that, for every fixed $0\leq\mu<1/2$, there is
an integer $q=q(\mu)>0$ such that homogeneous
$\Sigma\Pi\Sigma\Pi\Sigma$ circuits computing $\IMM_{d^q,d}$, with
$N=d^{2q+1}$ ambient variables and bottom fan-in at most $N^\mu$, require
$N^{\Omega_\mu(\sqrt d)}$ gates over every field
\cite[Theorem~1.1]{BeraChakrabarti2015}.
Their precise Theorem~4.2 proves the lower bound for a restricted IMM
polynomial, using the number of distinct variables feeding each bottom
linear gate as its support parameter
\cite[Theorem~4.2]{BeraChakrabarti2015}.
The Kumar--Saraf and Bera--Chakrabarti bounds use different width--degree
regimes from $w=d=n$, and the latter restricts the bottom linear forms.

For a different hard family, Kumar and Saptharishi proved, over each fixed
finite field $\mathbb F_q$,
$\exp(\Omega_q(\sqrt d))$ lower bounds for homogeneous depth-five circuits
without a bottom-support restriction
\cite[Theorem~1]{KumarSaptharishi2017}.
Their family is based on Nisan--Wigderson designs and lies in
$\mathrm{VNP}$.
Armand, Behera, and Tavenas have since extended this finite-field approach
to depth-five circuits whose top-product fan-in is bounded by a fixed
polynomial in the degree $d$, without homogeneity
\cite[Corollary~1.3]{ArmandBeheraTavenas2026}.
These finite-field theorems concern NW polynomials rather than IMM; both
the field hypothesis and the exponential rate differ from those in
Theorem~\ref{thm:main}.

Superpolynomial lower bounds for unrestricted-bottom depth-five circuits
over characteristic zero were already established by the constant-depth
results of Limaye, Srinivasan, and Tavenas
\cite{LimayeSrinivasanTavenas2021,LimayeSrinivasanTavenas2025}.
We use the theorem numbering of their ECCC revision~1.
Their product-depth convention counts multiplication gates on a path, so
our five-layer model has product-depth two.  In the low-degree regime
$d\leq(\log w)/100$, their Corollary~3 gives
$w^{\Omega(\sqrt d)}$ for homogeneous product-depth-two circuits over every
field~\cite[Corollary~3]{LimayeSrinivasanTavenas2021}.
Their general-circuit results also imply superpolynomial lower bounds for
$\IMM_{n,n}$ over $\C$
\cite[Corollary~4 and the following discussion]{LimayeSrinivasanTavenas2021}.
Bhargav, Dutta, and Saxena improved the constant-depth exponents in the
low-degree regime \cite{BhargavDuttaSaxena2022}.
Forbes subsequently established low-depth IMM lower bounds over arbitrary
fields in the regime $d=o(\log w)$~\cite[Theorem~10]{Forbes2024}.

Amireddy, Garg, Kayal, Saha, and Thankey gave a direct shifted-partials
approach.  For homogeneous product-depth-two formulas they obtain
$2^{-O(d)}w^{\Omega(\sqrt d)}$ when $w=\omega(d)$, over every field
\cite[Theorem~24]{AmireddyGargKayalSahaThankey2023}.
The degree-dependent loss and the hypothesis $w=\omega(d)$ prevent direct
substitution of $w=d=n$.  The comparison addressed here is quantitative:
we obtain the $\sqrt n$ exponent in the balanced regime $w=d=n$ for
homogeneous five-layer circuits with unrestricted bottom linear forms.

\subsection{Proof overview}

The proof uses a derivative-based linear-algebraic measure.  Dimensions of
partial-derivative spaces were used by Nisan and
Wigderson~\cite{NisanWigderson1995}; shifted partial derivatives were developed
in work of Kayal~\cite{Kayal2012} and Gupta, Kamath, Kayal, and
Saptharishi~\cite{GuptaKamathKayalSaptharishi2014}.
We define the operator needed here directly and prove its estimates.

Partition the matrix layers into two groups containing $k$ and $m$ layers,
where $k+m=n$ and $m-k$ is of order $\sqrt n$.  Let $V$ and $U$ be the
corresponding variable sets.  For a homogeneous polynomial $g$ of degree
$n$, let $g_{[k,m]}$ be the sum of its monomials having degree $k$ in $V$ and
degree $m$ in $U$.  Replace its $V$-variables by partial derivatives and
its $U$-variables by multiplication operators.  On polynomials of fixed
bidegree $(a,b)$ this gives a linear map
\[
 F_g=g_{[k,m]}(\partial_V,U)
\]
into polynomials of bidegree $(a-k,b+m)$.  We choose $a,b$ of order
$n^{5/2}$, write $D$ for the dimension of the source, and use
$R(g)=\rank F_g$ as the complexity measure.  Differentiation in $V$ commutes
with multiplication in $U$; the full substitution therefore respects
products, while the finite map is linear in $g$ and its rank is subadditive.
Section~\ref{sec:rank} gives the definitions and parameters.

We first bound the rank of a circuit in Section~\ref{sec:circuits}.
For a product of homogeneous factors, resolve each factor into its possible
$V$-degrees.  Apply first those factors whose $V$-degree exceeds their
proportional share $ke/n$ for a factor of degree $e$.  Their composition passes through
a smaller intermediate polynomial space, which bounds its rank.  The
distance of each proportional share from the integers controls the sum over
all such decompositions.  These are the same integer-degree discrepancies
that underlie the residue method of Amireddy, Garg, Kayal, Saha, and
Thankey \cite[Definition~2.1 and Lemma~3.1]{AmireddyGargKayalSahaThankey2022Full};
here they control intermediate homogeneous dimensions.
If all factor degrees are small, these distances
give the required saving.  If some factor has large degree, expand just one
such factor into its bottom products of linear forms.  This costs at most
one additional factor of the circuit size.  Proposition~\ref{prop:circuit}
gives, for a size-$S$ circuit,
\[
 R(g)\leq 2S^2D\exp(C\sqrt n)\,n^{-\sqrt n/100}
\]
with an absolute constant $C$.

The opposite estimate uses the paths in matrix multiplication.
Section~\ref{sec:moments} arranges the two groups of layers in a balanced
order, with no adjacent $V$-layers.  Declare the monomials
$z^M/\sqrt{M!}$ orthonormal, where $M$ is an exponent vector and
$M!$ is the product of its coordinate factorials.  In these bases,
differentiation and multiplication have coefficients given by square roots
of exponents and their successors.  For $g=\IMM_{n,n}$,
the first two trace moments of $F_g^*F_g$ become weighted counts of paths
and quadruples of paths.  Exact factorial moments for uniform weak
compositions compare these weights with independent geometric random
variables.  We verify the coefficient signs needed for that comparison.

At each layer, a contributing quadruple has one of two pairing patterns.
A change of pattern forces four vertex labels to agree and reduces the
number of choices by a factor of $n$.  Section~\ref{sec:path-sum} sums the
resulting weights, treating isolated corrections and the two endpoints
explicitly.  The balanced order controls the weight accumulated along each
run of one pattern.  The trace inequality
\[
 \rank H\geq\frac{(\tr H)^2}{\tr(H^2)}
 \qquad(H\ne0\text{ positive semidefinite})
\]
then yields $R(\IMM_{n,n})\geq D\exp(-C'\sqrt n)$, as stated in
Proposition~\ref{prop:imm}.  Section~\ref{sec:completion} compares the two
rank estimates and proves the field extension and matching upper bound.
